\section{问题定义与报告主线}

设输入语音为 $x_i$，标签为 $y_i \in \{0,1\}^4$，对应集合
\[
C=\{g,b,d,z\}.
\]
当前数据集是单标签四分类，每个样本恰好属于四类中的一类。为了便于和不同方法对接，代码仍保留一对多接口：先为每个标签输出分数 $s_{i,c}$，再由阈值或 $\arg\max$ 给出预测。

这份报告的顺序是固定的。先讲严格课内版，因为它最接近课堂工具，也最适合说明题目本身；再讲课外扩展版，因为它确实更强，但边界已经超出课堂；最后再把 AI 基线单独放到后面，不让它抢走主线。

需要先说明一点：当前 \path{build/mswc_course_gbdz_local_patch_v1} 目录保留的是整套实验结果，但其中 \texttt{course\_filterbank\_corr} 是课外扩展版，不是严格课内版。正文里如果提到“课堂方法”，指的是 \texttt{strict\_course\_fft}。
